\begin{formbox}
\textbf{Name:} Tianrui Ji \hfill \textbf{Student ID:} 2361367\\
\textbf{Role:} Project manager and algorithm implementation.

\textbf{1. My specific contributions to this project:}
\begin{enumerate}[leftmargin=2em,itemsep=0.25em,topsep=0.15em]
    \item I coordinated the AI-assisted development and integration of
    conservative FDM and two-dimensional P1 FEM in
    \path{project2/algorithm_section/src/discretisation.py}. The work covered
    shared face coefficients, interface resistance, triangle assembly and
    Dirichlet elimination. I prepared the algorithm-section text, figures
    and code for handover to the report-writing member.

    \item I organised the computational materials for convergence, interface
    accuracy, solver tolerance and cost comparisons. The solver interface in
    \path{project2/algorithm_section/src/solvers.py} uses SciPy's CG and sparse
    LU routines. I worked on explaining their roles and the distinction
    between algebraic, discretisation and physical flux errors. The separate
    independent validation work is documented by the testing member in
    \path{project2/test/VALIDATION_RUN_2026-11-06.md}.

    \item I organised the AI-assisted interface-PINN extension under
    \path{project2/I-PINNs/}, including training configuration, server runs
    and collection of checkpoints and logs. I requested that the revised
    presentation compare this method with FDM and FEM, and clarify joint
    training across the interface. I also coordinated the presentation
    revisions and prepared to explain the algorithm and extension sections.
\end{enumerate}

\textbf{2. The hardest conceptual obstacle I encountered:}
The main difficulty was connecting the implemented equations with the
different accuracy measures. In particular, a small linear residual does
not establish PDE accuracy, and agreement between two interface fluxes
does not establish agreement with the exact physical flux. Preparing the
presentation helped me separate these checks and explain what each measures.

\textbf{3. One piece of AI-generated output my team used, and what I changed:}
An initial AI-assisted extension slide mainly compared the interface network
with an ordinary PINN. I requested a comparison with our FDM and FEM work,
common evaluation points, and a clearer diagram of joint training. I also
asked for the archived device, optimisation budget and runtime to be
checked. The revised material reports the A100 recorded in the environment
file and avoids a cross-device speed ranking. These revisions are now included in the advanced-method section; see
the accompanying AI transparency contribution for the evaluation protocol.

\textbf{4. One thing I still do not fully understand:}
I would like to understand how to choose interface-loss weights
systematically and obtain an error guarantee near the interface. A small
training loss or good accuracy at sampled points does not by itself provide
such a guarantee.

\vfill
\textbf{Signature:} \underline{\hspace{4.5cm}} \hfill
\textbf{Date:} \underline{\hspace{2.5cm}}
\end{formbox}
